\chapter{La Geometria: l'Apparato
Macro}\label{la-geometria-lapparato-macro}

\begin{quote}
\textbf{Core argument:} ogni attore egemonico occupa una posizione
in uno spazio tridimensionale di stato (coesione, capacità,
legitimacy); la sua interazione strategica con altri attori follows un
payoff la cui forma è identica a quella del singolo coupling
discorsivo della Topografia; la stabilità del sistema egemonico è una
soglia \(T_s\) che misura il rapporto tra forze di coerenza e di
decoherence.
\end{quote}



\section{Perché un Chapter Sulla
Geometria}\label{perchuxe9-un-chapter-sulla-geometria}

Il Ch.~1 ha presentato la Topografia come micro-foundation comunicativa.
Questo chapter presenta l'apparato simmetrico al livello
macro-geopolitico: la \emph{Geometria dell'Egemonia} (Marcelino, 2026),
che applica formalism geometrico, teoria dei giochi e estensione
stocastica alla dinamica tra attori egemonici. Il Ch.~3 dimostrerà che
le due strutture sono \emph{unificabili} through l'operator di
Kuramoto. Per seguire quella proof, the reader ha bisogno
dell'apparato della Geometria --- questo chapter lo fornisce in forma
compatta.



\section{Il Problema Egemonico
Classico}\label{il-problema-egemonico-classico}

La tradizione critica ha trattato l'egemonia in registro qualitativo:
Gramsci (egemonia come direzione morale e culturale), Cox (egemonia come
articolazione di forze materiali, idee, istituzioni), Wallerstein (cicli
sistemici). Manca un \textbf{formalism unificato} che permetta
predizioni datate e falsificabili. La Geometria propone uno spazio
metrico delle posizioni egemoniche e una dinamica formal delle loro
interazioni.



\section{\texorpdfstring{Lo Spazio degli Stati Egemonici
\(\mathcal{H} = [0, 10]^3\)}{Lo Spazio degli Stati Egemonici \textbackslash mathcal\{H\} = {[}0, 10{]}\^{}3}}\label{lo-spazio-degli-stati-egemonici-mathcalh-0-103}

Ogni attore egemonico \(i\) è caratterizzato da uno stato
\(\mathbf{s}_i = (s_{i,1}, s_{i,2}, s_{i,3}) \in \mathcal{H}\):

\begin{itemize}
\tightlist
\item
  \(s_{i,1} = s_i\) = \textbf{internal cohesion} (integrazione politica,
  omogeneità decisionale, controllo del territorio);
\item
  \(s_{i,2} = s_{ii}\) = \textbf{material capacity} (PIL, popolazione,
  tecnologia, energia, forze armate);
\item
  \(s_{i,3} = s_{iii}\) = \textbf{ethical-discursive legitimacy} (soft
  power, legitimacy internazionale, qualità del discorso pubblico).
\end{itemize}

\Hypoth{} \textbf{Disambiguazione notezionale.} Lo spazio
\(\mathcal{H}\) degli stati egemonici è denoteto con la lettera \emph{H}
(per \emph{Hegemonic}) per evitare conflitto con \(\mathcal{S}\) del
Ch.~5, che indica la \emph{superficie critica}
\(T_s^{(2),\text{sys}} = 1\) in \(\mathbb{V}^{(4)}\). I due oggetti
hanno natura matematica e ruolo diversi: \(\mathcal{H}\) è uno spazio di
stato statico tridimensionale; \(\mathcal{S}\) è un'ipersurface
(sottovarietà di codimensione 1) nello spazio structural
quadridimensionale.

\Hypoth{} \textbf{Mappa con la Topografia.} Gli assi
\((s_i, s_{ii}, s_{iii})\) sono gli omologhi di scala degli assi
\((x, y, z)\) della Topografia. La calibrazione empirica di
\(\Sigma^{\text{geo}}\) starting from \(\Sigma^{\text{comm}}\) è discussa
nell'Appendix B della DEQ$^2$.

\subsection{La Qualità Strategica}\label{la-qualituxe0-strategica}

\[
Q_i = \frac{s_{i,1} + s_{i,2} + s_{i,3}}{\sqrt{3}}
\]

è la qualità complessiva dell'attore \(i\) --- formalmente identica al
\(Q_0\) della Topografia, applicata agli stati egemonici instead che alle
componenti del discorso.



\section{Il Payoff Egemonico e il Folk
Theorem}\label{il-payoff-egemonico-e-il-folk-theorem}

\subsection{Payoff di Periodo}\label{payoff-di-periodo}

L'attore \(i\) ottiene, per ogni periodo, un payoff:

\[
\Pi_i(t+1) = H_i \cdot Q_i \cdot (1 - \delta_i \cdot s_{i,3})
\]

where \(H_i\) è l'\textbf{impronta sistemica} (la ``massa egemonica'' ---
analogo macro di \(M_e\)) e \(\delta_i\) il discount factor. Il
termine \((1 - \delta_i s_{i,3})\) è il \textbf{erosion factor di
legitimacy}: la legitimacy conquistata si consuma nell'esercizio del
potere.

\Proved{} \textbf{Identità formal con Topografia (verificationta nel Ch.~3
\S{}3.1.3).} Il payoff egemonico è strutturalmente la stessa equation
della predictive function della Topografia
\(\Pi(t+1) = M_e Q_0 (1 - \delta_R z_R)\), con corrispondenza:

{\def\LTcaptype{none} % do not increment counter
{\small\begin{longtable}[]{@{}
  >{\raggedright\arraybackslash}p{(\linewidth - 2\tabcolsep) * \real{0.5000}}
  >{\raggedright\arraybackslash}p{(\linewidth - 2\tabcolsep) * \real{0.5000}}@{}}
\toprule\noalign{}
\begin{minipage}[b]{\linewidth}\raggedright
Topografia
\end{minipage} & \begin{minipage}[b]{\linewidth}\raggedright
Geometria
\end{minipage} \\
\midrule\noalign{}
\endhead
\bottomrule\noalign{}
\endlastfoot
\(M_e\) (communicative mass) & \(H_i\) (impronta sistemica) \\
\(Q_0\) (ethical quality) & \(Q_i\) (qualità strategica) \\
\(\delta_R\) (sconto del ricevente) & \(\delta_i\) (sconto
dell'attore) \\
\(z_R\) (complessità del ricevente) & \(s_{i,3}\) (legitimacy erosa
nell'uso) \\
\end{longtable}}
}

\subsection{Il Folk Theorem della
Geometria}\label{il-folk-theorem-della-geometria}

\Proved{} Per il gioco egemonico ripetuto, la cooperazione tra egemoni
(anziché conflitto totale) è sostenibile come Nash equilibrium
perfetto nei sottogiochi se e solo se:

\[
\delta_i \geq \delta^* = 0{,}50
\]

\Hypoth{} \textbf{Identità con il Discursive Folk Theorem della
Topografia.} La stessa \(\delta^* = 0{,}50\) del Folk Discorsivo per
l'Emittente. Questa coincidenza numerica è una firma matematica di
scale invariance --- la stessa struttura governa cooperazione
comunicativa micro e cooperazione egemonica macro. La DEQ$^2$ la dimostra
come theorem structural (Ch.~3).

\subsection{Implicazione Geopolitica del Folk
Theorem}\label{implicazione-geopolitica-del-folk-theorem}

\Hypoth{} Sistemi con \(\delta < 0{,}50\) --- basso time horizon,
alta preferenza per il presente --- non possono cooperare stabilmente.
Il populismo del \emph{short term} (cicli elettorali brevi, retorica
trimestrale dei mercati, \emph{engagement} algoritmico) è
strutturalmente incompatibile con la cooperazione egemonica. Predizione:
i sistemi politici che riducono \(\delta\) collettivo sotto \(0{,}50\)
entrano in regime di non-cooperazione structural, indipendingmente
dalla volontà degli attori.



\section{\texorpdfstring{La Soglia di Sgretolamento \(T_s\)
(DEQ$^1$)}{La Soglia di Sgretolamento T\_s (DEQ$^1$)}}\label{la-soglia-di-sgretolamento-t_s-dequxb9}

Il \emph{Memoriale Tecnico DEQ} (2026-04-23) follows, dalle equations del
payoff egemonico in regime stocastico, una soglia composita:

\[
T_s^{(1)} = \frac{s_{iii} \cdot z \cdot \delta}{\sigma \cdot \varepsilon_{\text{dis}}}
\]

where:

\begin{itemize}
\tightlist
\item
  \(s_{iii}\) = ethical-discursive legitimacy del sistema;
\item
  \(z\) = projection on the dialectical axis (eredità diretta della
  Topografia);
\item
  \(\delta\) = time horizon (Folk Theorem);
\item
  \(\sigma\) = stochastic volatility ambientale (estensione EDSD, \S{}2.5);
\item
  \(\varepsilon_{\text{dis}}\) = moral disengagement coefficient
  (Bandura 1999, Milgram 1974).
\end{itemize}

\Proved{} \textbf{Lettura.} Numeratore: forze di coerenza. Denominatore:
forze di decoherence. La critical threshold è \(T_s^{(1)} = 1\). Sopra,
sistema stabile; sotto, phase transition imminente.

\Conj{} \textbf{Limite della DEQ$^1$.} L'equation \emph{non distingue} tra
sistemi resilienti e antifragili, né tra attrito esogeno (volatilità) ed
endogeno (involution). Due sistemi con stesso \(T_s^{(1)}\) possono
avere traiettorie radicalmente diverse. Questa lacuna motiva
l'estensione DEQ$^2$ con \(A\) e \(\Omega\) (Ch.~3).



\section{L'Estensione Stocastica
EDSD}\label{lestensione-stocastica-edsd}

La \emph{Geometria dell'Egemonia} (Ch.~12-13) introduce
un'\textbf{Equazione Dinamica Stocastica del Decadimento} che modella
l'evoluzione temporale di ciascuno stato \(s_{i,k}\) come processo
stocastico:

\[
ds_{i,k} = \mu_k(\mathbf{s}_i, t)\, dt + \sigma_k(\mathbf{s}_i, t)\, dW_k(t)
\]

where \(W_k(t)\) è un processo di Wiener e \(\sigma_k\) è la volatilità
sull'asse \(k\). La \(\sigma\) che entra in \(T_s^{(1)}\) è una
\textbf{norma aggregata} del vettore \((\sigma_1, \sigma_2, \sigma_3)\),
in genere \(\sigma = \|\boldsymbol{\sigma}\|_2\).

\Hypoth{} \textbf{Confollowsnza per la DEQ$^2$.} L'antifragility \(A\) del
Ch.~3 è la \emph{convessità della risposta} a questa \(\sigma\) ---
collega EDSD alla formalizzazione di Taleb. Sistemi con \(A > 0\)
trasformano la varianza stocastica in opzionalità.



\section{Cosa la Geometria Porta alla
DEQ$^2$}\label{cosa-la-geometria-porta-alla-dequxb2}

{\def\LTcaptype{none} % do not increment counter
{\small\begin{longtable}[]{@{}
  >{\raggedright\arraybackslash}p{(\linewidth - 4\tabcolsep) * \real{0.3333}}
  >{\raggedright\arraybackslash}p{(\linewidth - 4\tabcolsep) * \real{0.3333}}
  >{\raggedright\arraybackslash}p{(\linewidth - 4\tabcolsep) * \real{0.3333}}@{}}
\toprule\noalign{}
\begin{minipage}[b]{\linewidth}\raggedright
Apparato Geometrico
\end{minipage} & \begin{minipage}[b]{\linewidth}\raggedright
Ruolo nella DEQ$^2$
\end{minipage} & \begin{minipage}[b]{\linewidth}\raggedright
Reference
\end{minipage} \\
\midrule\noalign{}
\endhead
\bottomrule\noalign{}
\endlastfoot
Spazio \(\mathcal{H} = [0,10]^3\) degli stati egemonici & Spazio degli
stati macro & Ch.~3 \\
\(Q_i\) qualità strategica & Forma identica a \(Q_0\) Topografia
(scale invariance) & Ch.~3 \S{}3.1.3 \\
\(\Pi_i(t+1) = H_i Q_i (1 - \delta_i s_{i,3})\) & Legge macro-payoff
identica a $\Pi$ Topografia & Ch.~3-4 \\
Folk Theorem \(\delta^* = 0{,}50\) & Stessa soglia di Folk Discorsivo
--- scale invariance & Ch.~3 \S{}3.1.3 \\
\(T_s^{(1)} = s_{iii} z \delta / (\sigma \varepsilon_{\text{dis}})\) &
Caso limite della DEQ$^2$ per \(A = \Omega = 0\) & Ch.~3 \\
EDSD: \(ds = \mu\, dt + \sigma\, dW\) & Definizione operativa di
\(\sigma\) in \(T_s\) & Ch.~4 \\
Coefficiente di disimpegno \(\varepsilon_{\text{dis}}\)
(Bandura/Milgram) & Variabile macro che assorbe la segnalazione
bayesiana di Topografia & Ch.~3 \\
\end{longtable}}
}

\Hypoth{} \textbf{Sintesi epistemica.} Se la Topografia è la
termodinamica statistica del sistema discorsivo, la Geometria è la
\textbf{termodinamica fenomenologica} del sistema egemonico --- descrive
le leggi di trasformazione macroscopiche senza esplicitarne la
micro-foundation. La DEQ$^2$ (Ch.~3-4) compie il passo che unifica i due
livelli, dimostrando che le leggi macro emergono per trasposizione di
Kuramoto dalle leggi micro.



\section{Closing Epistemic Notes
Chapter}\label{note-epistemiche-di-chiusura-chapter}

{\def\LTcaptype{none} % do not increment counter
{\small\begin{longtable}[]{@{}
  >{\raggedright\arraybackslash}p{(\linewidth - 6\tabcolsep) * \real{0.2500}}
  >{\raggedright\arraybackslash}p{(\linewidth - 6\tabcolsep) * \real{0.2500}}
  >{\raggedright\arraybackslash}p{(\linewidth - 6\tabcolsep) * \real{0.2500}}
  >{\raggedright\arraybackslash}p{(\linewidth - 6\tabcolsep) * \real{0.2500}}@{}}
\toprule\noalign{}
\begin{minipage}[b]{\linewidth}\raggedright
\#
\end{minipage} & \begin{minipage}[b]{\linewidth}\raggedright
Claim
\end{minipage} & \begin{minipage}[b]{\linewidth}\raggedright
Status
\end{minipage} & \begin{minipage}[b]{\linewidth}\raggedright
Reference
\end{minipage} \\
\midrule\noalign{}
\endhead
\bottomrule\noalign{}
\endlastfoot
1 & Spazio \(\mathcal{H}\) degli stati egemonici & \Proved{} definition
operativa & Geometria \S{}2 \\
2 & Forma di \(Q_i\) (identica a \(Q_0\) Topografia) & \Proved{}
verified & Geometria \S{}2 / DEQ$^2$ Ch.~3 \\
3 & Payoff \(\Pi_i\) identico a $\Pi$ Topografia & \Proved{} verified
(Ch.~3) & Geometria \S{}3 / DEQ$^2$ Ch.~3 \\
4 & Folk Theorem \(\delta^* = 0{,}50\) & \Proved{} followsto & Geometria
\S{}3 \\
5 & \(T_s^{(1)}\) come soglia composita & \Hypoth{} working hypothesis &
Memoriale DEQ \\
6 & EDSD come dinamica stocastica & \Proved{} formulazione standard
(Itô) & Geometria \S{}12-13 \\
7 & \(\varepsilon_{\text{dis}}\) ancorata a Bandura/Milgram & \Proved{}
ancoraggio empirico documentato & Memoriale DEQ \\
8 & Limite della DEQ$^1$ in the absence of \(A, \Omega\) & \Conj{} critica
costruttiva, motivante per DEQ$^2$ & DEQ$^2$ Ch.~3 \\
\end{longtable}}
}


